% =============================================================================
% 7. LIMITATIONS AND DISCUSSION
% =============================================================================
\section{Limitations and Discussion}
\label{sec:limitations}

We close by stating six limitations of the present study, both to scope
our claims and to flag concrete directions for follow-up work.

\paragraph{Limitation 1: Python-centric evaluation.}
The mid-training corpus is exclusively Python, and the in-domain agent
benchmarks (SWE-Bench-Verified and SWE-Bench-Lite) are likewise
Python-centric. Cross-language generalisation is supported only
indirectly, through the FullStackBench-ZH bilingual coding benchmark in
Section~\ref{sec:exp:gen}, and we do not cover Java, C++, Rust, or other
language ecosystems in the main tables. Whether the function-aware FIM
recipe transfers to languages with different module systems, type
disciplines, or call-resolution semantics is left to future work.

\paragraph{Limitation 2: Dependency on a frontier teacher for CoT.}
Our default recipe relies on the Gemini-3-Preview API to synthesise the
rationale embedded inside each FIM middle span. The CoT-source ablation
in Section~\ref{sec:exp:ablation} shows that self-generated rationales
recover most of the gain and FIM structure alone already beats the
no-mid-training baseline, so the recipe is not fundamentally
distillation-bound; nonetheless, a fully open-source replication of our
best configuration requires either a comparably strong open teacher or
acceptance of the residual gap reported in that ablation.

\paragraph{Limitation 3: Imperfect program dependency graph analysis on dynamic Python.}
Our PDG construction operates on the abstract syntax tree and resolves
calls via a combination of qualified-name matching and short-name
fallback. This works well for the vast majority of canonical Python code
but is imperfect on idioms involving dynamic dispatch, decorator chains,
and monkey patching, where the static call graph and the run-time call
graph diverge. We characterise these failure modes and the corresponding
mitigations in Appendix~\ref{app:algo}; an analysis-aware extension that
incorporates light dynamic information is a natural next step.

\paragraph{Limitation 4: Cross-base-family validation is partial.}
The cross-base-model evidence in Section~\ref{sec:exp:crossbase} comes
from a single non-Qwen2.5-Coder configuration (Qwen3-8B paired with the
swe-lego post-training pipeline). Because moving from Qwen2.5-Coder to
Qwen3-8B simultaneously varies the post-training pipeline (swe-lego is
required to fit Qwen3's $40{,}960$-token context), the cross-family
result should be read as ``the recipe is not specific to a single
pretraining/post-training combination'' rather than as a guarantee that
it generalises across all base-model families. Validations on Llama, the
DeepSeek family, and other architectures are not covered here.

\paragraph{Limitation 5: Capability-preservation comparison only at 14B.}
The controlled three-row comparison in Section~\ref{sec:exp:gen}
(instruct base vs.\ post-training only vs.\ mid-training plus
post-training) is conducted at the 14B scale only, due to the cost of
running the full set of non-agent coding and tool-use benchmarks at
every scale. The 7B and 32B in-domain SWE-Bench trends are consistent
with the 14B picture and suggest a similar regression-recovery pattern,
but we do not run the fully controlled triple at those scales and
therefore do not make a quantitative claim about capability preservation
outside 14B.

\paragraph{Limitation 6: Modularity assumption of function-aware selection.}
Function-aware FIM presupposes that the underlying code is meaningfully
modular, so that masking a single function or a small connected group
produces a non-trivial reasoning target. On extreme monolithic code or
inline-heavy codebases---scripts dominated by top-level statements,
generated code with few callable abstractions, or notebooks---the
selection pipeline has fewer eligible targets and may degrade
gracefully to a much smaller usable subset of files. We do not study
this regime systematically here; both characterising the breakdown point
and extending the algorithm with code-block-level abstractions for
non-modular code are open directions.
